\subsection*{Deriving \emph{mdspan}s}

New shapes often derive from existing ones. \textbf{Tiling} or \textbf{blocking} divides dimensions, while \textbf{padding} increases selected dimensions by a specified amount.

Shape arithmetic expresses these derivations directly:

\begin{lstlisting}[language=croqtile,caption={The original shape}]
shape : [128, 64]; // initial shape
new-shape0 : shape [(0) / 2, (1) / 4, 1];  // tile and reshape: [1, 64, 16]
new-shape1 : shape [(1) + 2, (0) / 16];    // pad and reshape: [66, 8]
\end{lstlisting}

The first derived shape divides dimension 0 by 2 and dimension 1 by 4, then adds a dimension of size 1. The expression determines the new dimension order. It changes the shape description; it does not move any elements.

In Croqtile, the definition of \texttt{new-shape0} is equivalent to:

\begin{lstlisting}[language=croqtile,caption={A derived shape}]
new-shape0: [shape(0) / 2, shape(1) / 4, 1];
\end{lstlisting}

This form lists the dimensions explicitly inside \texttt{[]}. The \textbf{element-of} operation \texttt{()} retrieves a dimension from the existing shape. Both forms produce the same result; the earlier form is syntactic sugar for the complete initialization expression.

The second derived shape adds 2 to the original dimension 1 and divides dimension 0 by 16, then places them in that order.

Furthermore, you may use the \texttt{mdspan} as a whole for derivations:

\begin{lstlisting}[language=croqtile,caption={Rebinding a shape}]
shape : [32, 72]
new-shape0 : shape;  // [32, 72]
new-shape1 : shape + 1;  // [33, 73]
new-shape2 : shape / 4;  // [8, 18]
new-shape3 : [shape, 6]; // [32, 72, 6]
\end{lstlisting}

Arithmetic on a whole shape applies to each dimension. Adding 1 to \co@[128, 64]@ therefore gives \co@[129, 65]@; it does not append a dimension. The two expressions below spell out that same calculation:

The expression \texttt{new-\allowbreak{}shape1 :\allowbreak{} shape + 1;}

is equivalent to

The expression \texttt{new-\allowbreak{}shape1 :\allowbreak{} shape [(0) + 1, (1) + 1];}

The derived definition of \texttt{new-shape3} demonstrates that using an \emph{mdspan} variable in an \emph{mdspan initialization expression} results in \textbf{concatenation} behavior. Thus, the declaration

The expression \texttt{new-\allowbreak{}shape3 :\allowbreak{} [shape, 6];}

is equivalent to

The expression \texttt{new-\allowbreak{}shape3 :\allowbreak{} shape [(0), (1), 6];}

The shorter form is syntactic sugar for the complete shape definition.
